% Generated draft: requires mathematical and source review.
\documentclass[11pt]{article}
\usepackage[margin=25mm]{geometry}
\usepackage{fontspec}
\setmainfont{Latin Modern Roman}
\usepackage{amsmath,amssymb,mathtools,unicode-math}
\setmathfont{Latin Modern Math}
\usepackage{enumitem,needspace,xurl,hyperref}
\setlength{\parindent}{0pt}
\setlength{\parskip}{5pt}
\setlength{\emergencystretch}{3em}
\newcommand{\N}{\mathbb{N}}
\newcommand{\Z}{\mathbb{Z}}
\newcommand{\Q}{\mathbb{Q}}
\newcommand{\R}{\mathbb{R}}
\newcommand{\C}{\mathbb{C}}
\newcommand{\F}{\mathbb{F}}
\newcommand{\A}{\mathbb{A}}
\newcommand{\PP}{\mathbb{P}}
\newcommand{\E}{\mathbb{E}}
\newcommand{\eps}{\varepsilon}
\newcommand{\dd}{\,\mathrm{d}}
\begin{document}
\section*{UnsolvedMath ID: 700916}
\textbf{Source problem code:} catalogue entry
\subsection{Definitions and mathematical statement}
A group \(G\) is \emph{periodic} if for every \(g\in G\) there exists an integer \(n\geq 1\) such that \(g^n=1_G\). No common bound on these orders is assumed. The group is \emph{locally finite} if \(\langle S\rangle\) is finite for every finite subset \(S\subseteq G\).

An automorphism \(\alpha\in\operatorname{Aut}(G)\) has order exactly \(4\) if \(\alpha^4=\operatorname{id}_G\) and \(\alpha^2\neq\operatorname{id}_G\). Interpret \emph{regular} in its usual fixed-point-free sense:
\[
\{g\in G:\alpha(g)=g\}=\{1_G\}.
\]
The supplied record does not define this term, so this interpretation requires confirmation against the original source. In particular, no fixed-point-free condition on \(\alpha^2\) is imposed here.

Under this interpretation, the question attributed to P. V. Shumyatsky is: for every periodic group \(G\), does the existence of a regular automorphism of order exactly \(4\) imply that
\[
\lvert\langle g_1,\ldots,g_r\rangle\rvert<\infty
\qquad\text{for every integer }r\geq 1\text{ and every }g_1,\ldots,g_r\in G?
\]
\subsection{Short English statement}
Must every periodic group admitting a fixed-point-free automorphism of order exactly \(4\) be locally finite? Here periodic means that every element has finite order, and locally finite means that every finitely generated subgroup is finite. The interpretation of ``regular'' as fixed-point-free requires source confirmation.
\subsection{Sources}
\begin{itemize}
\item Local UnsolvedMath registry, permanent ID \(700916\), exact source problem code \texttt{catalogue entry}, title \emph{Kourovka Notebook Problem 12.100}. The supplied registry labels the problem open; this status has not been independently verified.
\item \emph{The Kourovka Notebook (Unsolved Problems in Group Theory)}, original problem 12.100, 12th issue (1992); attributed to P. V. Shumyatsky. The supplied background also identifies the 21st issue edition (2026), arXiv:1401.0300v48. The collection metadata names E. I. Khukhro and V. D. Mazurov as editors.
\end{itemize}
\end{document}
